% FILE: chapters/ch04_s5_anomaly_detection.tex
\section{Anomaly Detection}\label{sec:anomaly_detection}

In the preceding sections of this chapter, we studied algorithms that classify observations into predefined categories or organize unlabeled data into natural geometric clusters. Yet across the physical world of the Internet of Things---from industrial gas turbines, chemical pipelines, and electrical power grids to embedded biometric wearables---the most critical operational question is rarely ``Which known category does this reading belong to?'' Instead, the urgent question is:
\begin{center}
\textbf{``Is this physical system behaving normally, or is something going dangerously wrong?''}
\end{center}

This is the domain of \textbf{Anomaly Detection} (also termed \textit{outlier detection}, \textit{novelty detection}, or \textit{fault identification}). An anomaly is an observation or pattern of observations that deviates so substantially from the established baseline of normal behavior that it arouses suspicion that it was generated by a different underlying physical mechanism, a degraded mechanical state, or an external cyber-physical intrusion.

\begin{figure}[htbp]
\centering
\begin{tikzpicture}[
    scale=0.92,
    box/.style={rectangle, draw=gray!70, fill=gray!5, thick, rounded corners=2mm, minimum height=1.0cm, align=center, font=\small\sffamily},
    anom_box/.style={rectangle, draw=warningred!80, fill=warningred!10, thick, rounded corners=2mm, minimum height=1.1cm, align=center, font=\small\sffamily},
    arrow/.style={-{Stealth[length=2.5mm]}, thick, draw=gray!80}
]
    % Central Input
    \node[box, minimum width=12.2cm, fill=keyideablue!10, draw=keyideablue!80] (input) at (0, 3.4) {\textbf{Streaming Telemetry Stream} $\mathbf{x}(t) \in \mathbb{R}^D$ (Vibration, Temperature, Current, Flow)};

    % Three Anomaly Categories
    \node[anom_box, minimum width=3.6cm] (point) at (-4.2, 1.4) {\textbf{1. Point Anomaly}\\ \footnotesize Instantaneous extreme spike;\\ \footnotesize e.g., Temp $= 145^\circ\text{C}$};
    \node[anom_box, minimum width=3.6cm] (context) at (0.0, 1.4) {\textbf{2. Contextual Anomaly}\\ \footnotesize Value normal overall, but\\ \footnotesize wrong for context (e.g. night)};
    \node[anom_box, minimum width=3.6cm] (collective) at (4.2, 1.4) {\textbf{3. Collective Anomaly}\\ \footnotesize Sequence individually valid,\\ \footnotesize but pattern indicates failure};

    % Arrows
    \draw[arrow] (-4.2, 2.9) -- (point.north);
    \draw[arrow] (0.0, 2.9) -- (context.north);
    \draw[arrow] (4.2, 2.9) -- (collective.north);

    % Diagnostic Action
    \node[font=\scriptsize\sffamily, text=gray!80] at (-4.2, 0.4) {Single-sample threshold violation};
    \node[font=\scriptsize\sffamily, text=gray!80] at (0.0, 0.4) {Conditional / state violation};
    \node[font=\scriptsize\sffamily, text=gray!80] at (4.2, 0.4) {Correlated trajectory violation};
\end{tikzpicture}
\caption{Taxonomy of operational anomalies in IoT and industrial systems. Detecting faults requires distinguishing between isolated measurement artifacts, context-dependent irregularities, and structural pattern breakdowns.}\label{fig:anomaly_types}
\end{figure}

Detecting anomalies in cyber-physical systems poses four formidable engineering challenges:
\begin{enumerate}
    \item \textbf{Extreme Class Imbalance}: Healthy operating data is overwhelmingly abundant (often exceeding $99.9\%$ of all telemetry), while mechanical faults, component fractures, and leaks are exceptionally rare.
    \item \textbf{Novelty of Unknown Failure Modes}: Machines rarely fail in identical ways. A newly installed centrifugal pump may develop cavitation, impeller unbalance, bearing spalling, or seal leakage. We cannot train a supervised classifier on classes that have never occurred.
    \item \textbf{Sensor Correlations}: Modern machinery is monitored by dozens of coupled sensors. A temperature of $42^\circ\text{C}$ and a pressure of $2.1\text{ bar}$ may both be perfectly valid operating values individually, but encountering high temperature simultaneously with minimum pressure may violate fundamental thermodynamic laws, signaling a severe system rupture.
    \item \textbf{Edge Resource Constraints}: Anomaly detection must frequently run locally on embedded microcontrollers (such as an ESP32 or ARM Cortex-M) mounted directly to the machine chassis to trigger emergency kill-switches within milliseconds without waiting for cloud round-trip latencies.
\end{enumerate}

In this section, we examine the four primary paradigms for anomaly detection:
\begin{itemize}
    \item \textbf{Statistical Thresholds}: Univariate Z-scores, Tukey's boxplot fences, and multivariate Mahalanobis distance.
    \item \textbf{Isolation Forest}: An efficient ensemble method based on random recursive space-partitioning.
    \item \textbf{One-Class Support Vector Machines (OC-SVM)}: Boundary estimation enclosing healthy data in a reproducing kernel Hilbert space.
    \item \textbf{Autoencoders}: Deep neural networks trained to compress and reconstruct normal patterns, leveraging reconstruction error as a quantitative anomaly score.
\end{itemize}

% ------------------------------------------------------------------------------
\subsection{Statistical Thresholds}\label{subsec:statistical_thresholds}
% ------------------------------------------------------------------------------

Statistical methods construct a probabilistic or geometric model of normal data. Observations that fall into regions of extremely low probability density or exceed defined statistical bounds are declared anomalous.

\subsubsection{Univariate Thresholds: Z-Score and Tukey's Fences}

When monitoring a single scalar sensor stream $x_1, x_2, \dots, x_N$, two classic statistical rules establish baseline operating thresholds:

\paragraph{The Gaussian Z-Score (3-Sigma Rule)}
Assuming the sensor telemetry is approximately normally distributed with empirical mean $\mu$ and standard deviation $\sigma$:
\begin{equation}
z_i = \frac{x_i - \mu}{\sigma}\label{eq:z_score}
\end{equation}
Under a Gaussian distribution $\mathcal{N}(\mu, \sigma^2)$, approximately $99.73\%$ of all observations fall within three standard deviations of the mean ($|z_i| \le 3$). A reading is flagged as an anomaly if:
\begin{equation}
|z_i| > \tau, \quad \text{typically } \tau = 3.0\label{eq:three_sigma_test}
\end{equation}
While computationally instantaneous (requiring only two floating-point operations per sample on a microcontroller), the 3-sigma test possesses two severe weaknesses:
\begin{itemize}
    \item \textbf{Contamination Sensitivity}: The sample mean $\mu = \frac{1}{N}\sum x_i$ and standard deviation $\sigma = \sqrt{\frac{1}{N}\sum (x_i - \mu)^2}$ are non-robust. A handful of massive outlier spikes during calibration will inflate $\mu$ and dramatically expand $\sigma$, masking subsequent anomalies.
    \item \textbf{Normality Assumption}: Physical sensor noise is frequently skewed, bimodal, or heavy-tailed, rendering theoretical Gaussian confidence levels invalid.
\end{itemize}

\paragraph{Tukey's Interquartile Range (IQR) Fences}
To eliminate sensitivity to extreme values and distributional assumptions, John Tukey formulated the non-parametric boxplot fence rule based on rank statistics:
\begin{enumerate}
    \item Sort the calibration dataset and calculate the first quartile $Q_1$ (25th percentile) and third quartile $Q_3$ (75th percentile).
    \item Compute the Interquartile Range ($\text{IQR}$), representing the spread of the central $50\%$ of the data:
    \begin{equation}
    \text{IQR} = Q_3 - Q_1\label{eq:iqr}
    \end{equation}
    \item Establish the inner fences:
    \begin{equation}
    F_{\text{lower}} = Q_1 - 1.5 \times \text{IQR}, \qquad F_{\text{upper}} = Q_3 + 1.5 \times \text{IQR}\label{eq:tukey_fences}
    \end{equation}
\end{enumerate}
Any sensor reading falling outside $[F_{\text{lower}}, F_{\text{upper}}]$ is classified as an outlier. For extreme outlier detection, an expansion factor of $3.0 \times \text{IQR}$ is utilized (outer fences). Because $Q_1$, $Q_3$, and $\text{IQR}$ depend on order statistics rather than squared distances, Tukey's fences tolerate up to $25\%$ arbitrary contamination without breakdown.

\subsubsection{The Failure of Univariate Thresholds in Coupled Systems}

Applying independent univariate thresholds to multidimensional telemetry produces an axis-aligned bounding box (or hyper-rectangle in higher dimensions). Figure~\ref{fig:univariate_vs_multivariate} illustrates why this strategy fails catastrophically when physical sensors are correlated.

\begin{figure}[htbp]
\centering
\begin{tikzpicture}[scale=0.95]
    % Axes
    \draw[-{Stealth[length=2.5mm]}, thick, gray!80] (0, 0) -- (7.5, 0) node[right, font=\small\sffamily] {Pump Temperature ($x_1$)};
    \draw[-{Stealth[length=2.5mm]}, thick, gray!80] (0, 0) -- (0, 6.2) node[above, font=\small\sffamily] {Discharge Pressure ($x_2$)};

    % Univariate Box
    \draw[dashed, warningred!70, fill=warningred!5, thick] (1.2, 0.8) rectangle (6.2, 5.2);
    \node[warningred!80, font=\scriptsize\sffamily, anchor=south west] at (1.2, 5.2) {Univariate Bounding Box ($3\sigma_1 \times 3\sigma_2$)};

    % Ellipse for Correlated Normal
    \draw[keyideablue!80, fill=keyideablue!12, thick, rotate around={38:(3.7, 3.0)}] (3.7, 3.0) ellipse (2.6cm and 0.85cm);
    \node[keyideablue!90, font=\scriptsize\sffamily] at (3.7, 3.0) {\textbf{Healthy Physical Envelope}};

    % Normal Data Points
    \foreach \x/\y in {2.2/1.8, 2.6/2.2, 3.0/2.3, 3.2/2.7, 3.6/3.1, 3.8/2.9, 4.0/3.4, 4.4/3.5, 4.8/4.0, 5.1/4.2} {
        \fill[keyideablue!80] (\x, \y) circle (2pt);
    }

    % Correlation-Breaking Anomaly (Inside box, outside ellipse)
    \coordinate (anom1) at (2.0, 4.6);
    \fill[warningred] (anom1) circle (3pt);
    \draw[warningred, thick] (anom1) circle (6pt);
    \draw[-{Stealth[length=2mm]}, warningred, thick] (3.0, 5.5) -- (anom1);
    \node[warningred, font=\scriptsize\sffamily, align=center, anchor=south] at (3.0, 5.5) {\textbf{Correlation-Breaking Anomaly}\\Low Temp + Max Pressure\\(\textit{Inside univariate box!})};

    % Extreme Point Anomaly (Outside both)
    \coordinate (anom2) at (7.0, 5.8);
    \fill[warningred] (anom2) circle (3pt);
    \draw[warningred, thick] (anom2) circle (6pt);
    \node[warningred, font=\scriptsize\sffamily, anchor=west] at (7.1, 5.8) {\textbf{Point Spike}\\(Outside both)};
\end{tikzpicture}
\caption{The failure of univariate bounds on correlated variables. The red circle at $(2.0, 4.6)$ represents a cold pump operating at near-bursting discharge pressure. Univariate thresholds miss this severe fault completely because both values fall within their individual $3\sigma$ ranges. The Mahalanobis distance detects it immediately by accounting for covariance.}\label{fig:univariate_vs_multivariate}
\end{figure}

\subsubsection{Multivariate Anomaly Detection: Mahalanobis Distance}

To capture covariance across multiple sensor channels, we employ the \textbf{Mahalanobis Distance}. Let $\mathbf{x} = [x_1, x_2, \dots, x_D]^T \in \mathbb{R}^D$ represent a multivariate sensor reading. Let $\boldsymbol{\mu} \in \mathbb{R}^D$ be the empirical mean vector and $\mathbf{\Sigma} \in \mathbb{R}^{D \times D}$ be the covariance matrix estimated from healthy operational baseline data:
\begin{equation}
\boldsymbol{\mu} = \frac{1}{N}\sum_{i=1}^N \mathbf{x}_i, \qquad \mathbf{\Sigma} = \frac{1}{N-1}\sum_{i=1}^N (\mathbf{x}_i - \boldsymbol{\mu})(\mathbf{x}_i - \boldsymbol{\mu})^T\label{eq:sample_mean_cov}
\end{equation}
The squared Mahalanobis distance from observation $\mathbf{x}$ to the normal cluster center $\boldsymbol{\mu}$ is defined as:
\begin{equation}
D_M^2(\mathbf{x}) = (\mathbf{x} - \boldsymbol{\mu})^T \mathbf{\Sigma}^{-1} (\mathbf{x} - \boldsymbol{\mu})\label{eq:mahalanobis_dist}
\end{equation}
Notice how the covariance matrix inverse $\mathbf{\Sigma}^{-1}$ de-correlates the features and scales each direction inversely proportional to its variance. Geometrically, level sets of $D_M^2(\mathbf{x}) = c$ form concentric ellipsoids oriented along the principal axes of data variation.

\begin{keyidea}[Chi-Squared Hypothesis Testing for Anomaly Thresholds]
If the healthy baseline data follows a $D$-dimensional multivariate Gaussian distribution $\mathcal{N}(\boldsymbol{\mu}, \mathbf{\Sigma})$, the squared Mahalanobis distance follows a Chi-squared distribution with $D$ degrees of freedom:
\begin{equation}
D_M^2(\mathbf{x}) \sim \chi^2_D\label{eq:chi_sq_distribution}
\end{equation}
To detect anomalies at a statistical significance level of $\alpha$ (e.g., $\alpha = 0.01$ for a $99\%$ confidence interval):
\begin{equation}
\text{Flag } \mathbf{x} \text{ as an anomaly if } D_M^2(\mathbf{x}) > \chi^2_{D, 1-\alpha}\label{eq:mahalanobis_threshold}
\end{equation}
For example, in a 3-sensor system ($D = 3$) at a $99\%$ confidence level ($1 - \alpha = 0.99$), the critical value from the Chi-squared distribution table is $\chi^2_{3, 0.99} \approx 11.345$. Any sample with $D_M^2(\mathbf{x}) > 11.345$ ($D_M > 3.368$) is flagged as anomalous with only a $1\%$ theoretical false alarm rate under nominal conditions.
\end{keyidea}

\begin{example}[Calculating Mahalanobis Distance for a Bearing Sensor Pair]
A bearing monitoring module tracks temperature $x_1$ ($^\circ\text{C}$) and vibration $x_2$ ($\text{mm/s}$). Baseline healthy operation yields:
\begin{equation*}
\boldsymbol{\mu} = \begin{bmatrix} 30.0 \\ 0.20 \end{bmatrix}, \qquad \mathbf{\Sigma} = \begin{bmatrix} 4.0 & 0.04 \\ 0.04 & 0.0016 \end{bmatrix}
\end{equation*}
A new telemetry reading arrives: $\mathbf{x} = [32.0, 0.24]^T$.
\begin{enumerate}
    \item Compute difference vector $\mathbf{d} = \mathbf{x} - \boldsymbol{\mu} = [32.0 - 30.0, \ 0.24 - 0.20]^T = [2.0, \ 0.04]^T$.
    \item Compute determinant $\det(\mathbf{\Sigma}) = (4.0)(0.0016) - (0.04)^2 = 0.0064 - 0.0016 = 0.0048$.
    \item Invert the $2 \times 2$ covariance matrix:
    \begin{equation*}
    \mathbf{\Sigma}^{-1} = \frac{1}{0.0048} \begin{bmatrix} 0.0016 & -0.04 \\ -0.04 & 4.0 \end{bmatrix} = \begin{bmatrix} 0.3333 & -8.3333 \\ -8.3333 & 833.333 \end{bmatrix}
    \end{equation*}
    \item Multiply $\mathbf{\Sigma}^{-1} \mathbf{d}$:
    \begin{equation*}
    \mathbf{\Sigma}^{-1} \mathbf{d} = \begin{bmatrix} 0.3333(2.0) - 8.3333(0.04) \\ -8.3333(2.0) + 833.333(0.04) \end{bmatrix} = \begin{bmatrix} 0.6667 - 0.3333 \\ -16.6667 + 33.3333 \end{bmatrix} = \begin{bmatrix} 0.3333 \\ 16.6667 \end{bmatrix}
    \end{equation*}
    \item Compute squared distance:
    \begin{equation*}
    D_M^2(\mathbf{x}) = \mathbf{d}^T (\mathbf{\Sigma}^{-1} \mathbf{d}) = 2.0(0.3333) + 0.04(16.6667) = 0.6667 + 0.6667 = \mathbf{1.3333}
    \end{equation*}
    $D_M(\mathbf{x}) = \sqrt{1.3333} \approx 1.1547$. For $D = 2$ at $\alpha = 0.05$, $\chi^2_{2, 0.95} = 5.991$. Because $1.3333 < 5.991$, the sample is safely inside the healthy envelope.
\end{enumerate}
\end{example}

% ------------------------------------------------------------------------------
\subsection{Isolation Forest}\label{subsec:isolation_forest}
% ------------------------------------------------------------------------------

Traditional anomaly detection algorithms profile normal instances and identify anomalies as samples that do not conform to that profile. In 2008, Fei Tony Liu, Kai Ming Ting, and Zhi-Hua Zhou introduced a radically different paradigm: \textbf{Isolation Forest} (\texttt{iForest}).

\begin{analogy}[Finding an Outlier in a Dense Crowd]
Imagine a crowd of $1{,}000$ people gathered closely together in the center of a football field, while a solitary individual stands completely alone near the corner flag. 

If you place random vertical and horizontal dividing ropes across the field at random coordinates:
\begin{itemize}
    \item The person standing alone near the corner flag will be isolated into their own private rectangular cell after only one or two random ropes.
    \item In contrast, isolating any specific individual standing inside the tightly packed crowd in the center requires dozens of crisscrossing ropes slicing through the swarm.
\end{itemize}
Anomalies are ``few and different.'' Because they reside in sparse regions of the feature space, they require far fewer random splits to isolate than normal points.
\end{analogy}

\begin{figure}[htbp]
\centering
\begin{tikzpicture}[scale=0.92]
    % Left: Partitioning diagram
    \draw[thick, gray!80] (0, 0) rectangle (5.5, 5.0);
    \node[font=\footnotesize\sffamily, gray!80] at (2.75, 5.25) {Feature Space Partitioning};

    % Dense cluster of normal points
    \foreach \x/\y in {3.2/2.1, 3.4/2.5, 3.8/2.2, 4.0/2.7, 4.2/2.3, 3.6/3.0, 3.9/3.2, 4.5/2.8, 3.5/2.6, 4.1/2.1} {
        \fill[keyideablue!70] (\x, \y) circle (2pt);
    }
    \node[keyideablue!90, font=\scriptsize\sffamily] at (3.9, 1.7) {Normal Cluster};

    % Isolated anomaly
    \coordinate (anom) at (0.8, 4.2);
    \fill[warningred] (anom) circle (3pt);
    \node[warningred, font=\scriptsize\sffamily, anchor=west] at (1.0, 4.2) {$\mathbf{x}_{\text{anom}}$};

    % Splits isolating anomaly
    \draw[warningred, thick, dashed] (0, 3.6) -- (5.5, 3.6) node[right, font=\tiny\sffamily] {Cut 1};
    \draw[warningred, thick, dashed] (1.8, 3.6) -- (1.8, 5.0) node[above, font=\tiny\sffamily] {Cut 2};
    \node[warningred, font=\tiny\bfseries\sffamily] at (0.9, 4.7) {Isolated! ($h=2$)};

    % Splits inside dense cluster
    \draw[gray!60, thin, dotted] (3.0, 0) -- (3.0, 3.6);
    \draw[gray!60, thin, dotted] (3.0, 2.0) -- (5.5, 2.0);
    \draw[gray!60, thin, dotted] (3.7, 2.0) -- (3.7, 3.6);
    \draw[gray!60, thin, dotted] (3.7, 2.8) -- (5.5, 2.8);
    \node[font=\tiny\sffamily, text=gray!80] at (4.6, 3.4) {Many cuts ($h=8$)};

    % Right: Tree representation
    \begin{scope}[shift={(7.2, 2.4)}]
        \node[font=\footnotesize\sffamily, gray!80] at (2.0, 2.85) {Isolation Tree (\textit{iTree})};

        % Tree nodes
        \node[circle, draw=gray!80, fill=gray!10, inner sep=2pt, font=\tiny\sffamily] (r) at (2.0, 2.2) {Root};
        \node[circle, draw=warningred, fill=warningred!20, inner sep=2pt, font=\tiny\sffamily] (n1) at (0.8, 1.1) {Split 1};
        \node[circle, draw=gray!80, fill=gray!10, inner sep=2pt, font=\tiny\sffamily] (n2) at (3.2, 1.1) {Normal Subtree};
        
        \node[rectangle, draw=warningred, fill=warningred!30, inner sep=2pt, font=\tiny\bfseries\sffamily] (leaf_anom) at (0.2, 0.0) {$\mathbf{x}_{\text{anom}}$};
        \node[circle, draw=gray!60, fill=gray!5, inner sep=1.5pt, font=\tiny\sffamily] (n3) at (1.4, 0.0) {$\dots$};

        \node[circle, draw=gray!60, fill=gray!5, inner sep=1.5pt, font=\tiny\sffamily] (d1) at (2.6, 0.0) {$\dots$};
        \node[circle, draw=gray!60, fill=gray!5, inner sep=1.5pt, font=\tiny\sffamily] (d2) at (3.8, 0.0) {$\dots$};

        \draw[thick, gray!70] (r) -- (n1);
        \draw[thick, gray!70] (r) -- (n2);
        \draw[thick, warningred] (n1) -- (leaf_anom);
        \draw[thick, gray!50] (n1) -- (n3);
        \draw[thick, gray!50] (n2) -- (d1);
        \draw[thick, gray!50] (n2) -- (d2);

        \draw[thick, -{Stealth[length=2mm]}, warningred] (-0.4, 1.8) -- node[left, font=\tiny\bfseries\sffamily, warningred] {Short Path ($h=2$)} (-0.4, 0.1);
    \end{scope}
\end{tikzpicture}
\caption{The mechanics of Isolation Forest. Left: Random axis-aligned cuts isolate an anomalous point in only two partitions, whereas points residing in the dense cluster require numerous recursive divisions. Right: In the resulting binary tree (\textit{iTree}), anomalous points terminate in shallow external leaves with small path length $h$.}\label{fig:isolation_forest_concept}
\end{figure}

\subsubsection{Algorithm: Building an Isolation Tree (\textit{iTree})}

An Isolation Tree is a proper binary tree constructed recursively as follows:
\begin{enumerate}
    \item Given a dataset $\mathcal{X} \subset \mathbb{R}^D$ of size $n$, if $n \le 1$ or tree depth reaches the limit $h_{\max} = \lceil \log_2(\psi) \rceil$, return an external leaf node.
    \item Randomly select a feature coordinate $q \in \{1, 2, \dots, D\}$ with uniform probability.
    \item Randomly select a split point $p$ uniformly at random between the minimum and maximum values of feature $q$ in the current subset:
    \begin{equation}
    p \sim \text{Uniform}(\min(\mathcal{X}_q), \ \max(\mathcal{X}_q))\label{eq:random_split}
    \end{equation}
    \item Partition $\mathcal{X}$ into left subset $\mathcal{X}_L = \{\mathbf{x} \in \mathcal{X} \mid x_q < p\}$ and right subset $\mathcal{X}_R = \{\mathbf{x} \in \mathcal{X} \mid x_q \ge p\}$.
    \item Recursively build the left child from $\mathcal{X}_L$ and right child from $\mathcal{X}_R$.
\end{enumerate}

An \textbf{Isolation Forest} builds an ensemble of $t$ independent trees (typically $t = 100$), each trained on a small randomly drawn sub-sample of size $\psi$ without replacement (typically $\psi = 256$).

\subsubsection{The Mathematical Anomaly Score}

To evaluate a query vector $\mathbf{x}$, it is traversed through all $t$ trees. For each tree $k$, the path length $h_k(\mathbf{x})$ is the number of edges traversed from the root to the terminating leaf node.

Because an \textit{iTree} is mathematically equivalent to a Binary Search Tree (BST), the average path length of an unsuccessful search among $n$ instances serves as the theoretical normalizer:
\begin{equation}
c(n) = 2 \ln(n - 1) + 2\gamma - \frac{2(n - 1)}{n}\label{eq:c_factor}
\end{equation}
where $\gamma \approx 0.5772156649$ is the Euler-Mascheroni constant. The term $c(n)$ represents the expected tree depth for a dataset of size $n$.

The \textbf{Anomaly Score} $s(\mathbf{x}, n)$ over an ensemble of $t$ trees is defined as:
\begin{equation}
s(\mathbf{x}, n) = 2^{-\frac{\mathbb{E}[h(\mathbf{x})]}{c(n)}}\label{eq:iforest_score}
\end{equation}
where $\mathbb{E}[h(\mathbf{x})] = \frac{1}{t}\sum_{k=1}^t h_k(\mathbf{x})$ is the mean path length across the forest.

The anomaly score $s \in (0, 1)$ possesses an intuitive mathematical interpretation:
\begin{itemize}
    \item \textbf{Definite Anomaly} ($s \to 1$): When $\mathbb{E}[h(\mathbf{x})] \to 0$, $s \to 2^0 = 1$. The sample terminates at very shallow leaf nodes across all trees, isolating almost immediately.
    \item \textbf{Normal Instance} ($s < 0.5$): When $\mathbb{E}[h(\mathbf{x})] \to n - 1$, $s \to 0$. The instance is deeply nested within dense clusters, requiring many splits.
    \item \textbf{Indeterminate Distribution} ($s \approx 0.5$): When $\mathbb{E}[h(\mathbf{x})] \approx c(n)$, $s \approx 2^{-1} = 0.5$. If all samples in an entire dataset return $s \approx 0.5$, the data has no distinctly isolated outliers.
\end{itemize}
In production, a common threshold is $s(\mathbf{x}) \ge 0.6$ or selecting the top $k$-th percentile corresponding to expected contamination.

\begin{keyidea}[Why Subsampling ($\psi = 256$) Eliminates Swamping and Masking]
Unlike decision trees in supervised learning that demand large training sets, Isolation Forest performs best when trained on small subsets ($\psi = 256$). Subsampling mitigates two major failure modes:
\begin{itemize}
    \item \textbf{Swamping}: When normal instances reside near outliers, large datasets cause normal instances to be grouped with outliers and mistakenly isolated early.
    \item \textbf{Masking}: When multiple outliers cluster together, they mimic a normal dense group, requiring many cuts to separate. A sub-sample of $256$ points drastically reduces the chance of selecting multiple correlated outliers, isolating the anomaly cleanly.
\end{itemize}
Subsampling also makes Isolation Forest exceptionally fast: training $100$ trees takes $\mathcal{O}(t \cdot \psi \log \psi)$, requiring merely milliseconds.
\end{keyidea}

% ------------------------------------------------------------------------------
\subsection{One-Class Support Vector Machines (OC-SVM)}\label{subsec:one_class_svm}
% ------------------------------------------------------------------------------

In Section~\ref{sec:classification}, we examined Support Vector Machines for distinguishing between two labeled classes by maximizing the separating margin. In the unsupervised setting, Bernhard Schölkopf et al.~extended this geometric framework into the \textbf{One-Class SVM} (OC-SVM).

\subsubsection{Geometric Formulation: Separating Normal Data from the Origin}

The fundamental premise of Schölkopf's One-Class SVM is to map training feature vectors $\mathbf{x}_i \in \mathbb{R}^D$ into a high-dimensional reproducing kernel Hilbert space $\mathcal{H}$ via a non-linear feature map $\phi(\mathbf{x})$. In that feature space, the algorithm finds the maximum margin hyperplane that separates the mapped training points from the coordinate origin $\mathbf{0}$:
\begin{equation}
\min_{\mathbf{w}, \boldsymbol{\xi}, \rho} \quad \frac{1}{2} \|\mathbf{w}\|^2 + \frac{1}{\nu N}\sum_{i=1}^N \xi_i - \rho\label{eq:ocsvm_primal}
\end{equation}
subject to the linear margin constraints:
\begin{equation}
\mathbf{w}^T \phi(\mathbf{x}_i) \ge \rho - \xi_i, \qquad \xi_i \ge 0 \quad \forall i \in \{1, \dots, N\}\label{eq:ocsvm_constraints}
\end{equation}
where:
\begin{itemize}
    \item $\mathbf{w}$ is the normal vector to the separating hyperplane in $\mathcal{H}$.
    \item $\rho$ is the orthogonal offset of the hyperplane from the origin.
    \item $\xi_i$ are non-negative slack variables permitting a fraction of training points to fall on the wrong side of the margin (handling baseline noise).
    \item $\nu \in (0, 1]$ is a critical user-specified hyperparameter.
\end{itemize}

\begin{keyidea}[The Physical Meaning of the Hyperparameter $\nu$]
The parameter $\nu \in (0, 1]$ fulfills two complementary mathematical roles:
\begin{enumerate}
    \item An \textbf{upper bound} on the fraction of training errors (outliers/false alarms) allowed inside the training baseline.
    \item A \textbf{lower bound} on the fraction of training samples that become support vectors.
\end{enumerate}
Setting $\nu = 0.05$ instructs the optimization to treat at most $5\%$ of the training data as potential outliers, tightly fitting a boundary around the remaining $95\%$.
\end{keyidea}

\subsubsection{The Kernel Trick and Decision Function}

By solving the Lagrangian dual problem, the optimal normal vector expands as a linear combination of support vectors: $\mathbf{w} = \sum_{i=1}^N \alpha_i \phi(\mathbf{x}_i)$. Using the kernel trick $K(\mathbf{x}_i, \mathbf{x}_j) = \langle \phi(\mathbf{x}_i), \phi(\mathbf{x}_j) \rangle$, the dual optimization becomes:
\begin{equation}
\min_{\boldsymbol{\alpha}} \quad \frac{1}{2} \sum_{i=1}^N \sum_{j=1}^N \alpha_i \alpha_j K(\mathbf{x}_i, \mathbf{x}_j) \quad \text{s.t.} \quad 0 \le \alpha_i \le \frac{1}{\nu N}, \quad \sum_{i=1}^N \alpha_i = 1\label{eq:ocsvm_dual}
\end{equation}
For industrial telemetry, the standard kernel choice is the Radial Basis Function (RBF / Gaussian) kernel:
\begin{equation}
K(\mathbf{x}, \mathbf{z}) = \exp(-\gamma \|\mathbf{x} - \mathbf{z}\|^2)\label{eq:rbf_kernel_ocsvm}
\end{equation}
Once $\boldsymbol{\alpha}$ and the threshold offset $\rho$ are solved, the decision function for a new sample $\mathbf{x}$ is:
\begin{equation}
f(\mathbf{x}) = \text{sign}\left( \sum_{i=1}^N \alpha_i K(\mathbf{x}_i, \mathbf{x}) - \rho \right)\label{eq:ocsvm_decision}
\end{equation}
If $f(\mathbf{x}) = +1$, the observation lies inside the learned non-linear envelope (nominal operation). If $f(\mathbf{x}) = -1$, the observation is flagged as an anomaly.

\begin{warning}[The Extreme Sensitivity of OC-SVM to the Kernel Scale $\gamma$]
While OC-SVM creates remarkably flexible non-linear decision boundaries, it is hypersensitive to the RBF parameter $\gamma$:
\begin{itemize}
    \item \textbf{$\gamma$ too small} (oversmoothing): The RBF kernel acts like a linear kernel. The boundary becomes an uninformative wide sphere enclosing vast regions of empty feature space, causing massive false negatives.
    \item \textbf{$\gamma$ too large} (overfitting): The Gaussian kernels decay to zero almost immediately around each individual training point. The decision boundary fragments into tiny isolated spheres centered on individual training vectors, causing rampant false alarms on normal operational variations.
\end{itemize}
Because ground-truth anomaly labels are unavailable during unsupervised training, tuning $\gamma$ requires heuristic grid search over the median pairwise distance ($1 / (2\sigma^2)$) or cross-validation on synthetic noise.
\end{warning}

% ------------------------------------------------------------------------------
\subsection{Autoencoders for Semi-Supervised Anomaly Detection}\label{subsec:autoencoders}
% ------------------------------------------------------------------------------

As sensor systems scale to high dimensions (such as 3-axis high-speed accelerometers recording full vibration frequency spectra with $D = 256$ to $1{,}024$ frequency bins), distance metrics and kernel evaluations suffer from the curse of dimensionality. Here, deep neural \textbf{Autoencoders} provide an exceptionally powerful representation-learning approach.

\subsubsection{Architecture: The Information Bottleneck}

An autoencoder is a feedforward neural network trained to execute an identity mapping: reproducing its input at its output layer ($\hat{\mathbf{x}} \approx \mathbf{x}$). Crucially, the network is constrained by a central \textbf{latent bottleneck} layer whose dimensionality $d$ is substantially smaller than the input dimension $D$ ($d \ll D$):

\begin{figure}[htbp]
\centering
\begin{tikzpicture}[
    scale=0.9,
    neuron/.style={circle, draw=keyideablue!80, fill=keyideablue!15, thick, minimum size=6.5mm, inner sep=0pt, font=\scriptsize\sffamily},
    latent/.style={circle, draw=exampleorange!80, fill=exampleorange!25, thick, minimum size=6.5mm, inner sep=0pt, font=\scriptsize\sffamily},
    out_neuron/.style={circle, draw=analogygreen!80, fill=analogygreen!15, thick, minimum size=6.5mm, inner sep=0pt, font=\scriptsize\sffamily},
    layer_label/.style={font=\small\bfseries\sffamily, align=center}
]
    % Input Layer (D = 5)
    \foreach \i in {1,...,5} {
        \node[neuron] (in\i) at (0, 3.6 - 0.9*\i) {$x_\i$};
    }
    \node[layer_label, keyideablue] at (0, 3.3) {Input $\mathbf{x} \in \mathbb{R}^D$};

    % Hidden Encoder Layer (4)
    \foreach \i in {1,...,4} {
        \node[neuron, fill=keyideablue!8] (enc\i) at (2.4, 3.15 - 0.9*\i) {};
    }
    \node[layer_label, text=gray!80] at (2.4, 3.3) {Encoder};

    % Bottleneck Latent Layer (d = 2)
    \foreach \i in {1,...,2} {
        \node[latent] (lat\i) at (4.8, 2.25 - 0.9*\i) {$z_\i$};
    }
    \node[layer_label, exampleorange] at (4.8, 3.3) {Bottleneck $\mathbf{z} \in \mathbb{R}^d$};

    % Hidden Decoder Layer (4)
    \foreach \i in {1,...,4} {
        \node[neuron, fill=analogygreen!8] (dec\i) at (7.2, 3.15 - 0.9*\i) {};
    }
    \node[layer_label, text=gray!80] at (7.2, 3.3) {Decoder};

    % Output Reconstruction Layer (5)
    \foreach \i in {1,...,5} {
        \node[out_neuron] (out\i) at (9.6, 3.6 - 0.9*\i) {$\hat{x}_\i$};
    }
    \node[layer_label, analogygreen] at (9.6, 3.3) {Recon $\hat{\mathbf{x}} \in \mathbb{R}^D$};

    % Connect Input to Enc
    \foreach \i in {1,...,5} {
        \foreach \j in {1,...,4} {
            \draw[gray!30, thin] (in\i) -- (enc\j);
        }
    }
    % Connect Enc to Latent
    \foreach \i in {1,...,4} {
        \foreach \j in {1,...,2} {
            \draw[gray!40, thin] (enc\i) -- (lat\j);
        }
    }
    % Connect Latent to Dec
    \foreach \i in {1,...,2} {
        \foreach \j in {1,...,4} {
            \draw[gray!40, thin] (lat\i) -- (dec\j);
        }
    }
    % Connect Dec to Out
    \foreach \i in {1,...,4} {
        \foreach \j in {1,...,5} {
            \draw[gray!30, thin] (dec\i) -- (out\j);
        }
    }

    % Compression Annotation
    \draw[thick, -{Stealth[length=2.5mm]}, keyideablue] (0.8, -1.8) -- node[below, font=\scriptsize\sffamily] {Compression: $f_\theta(\mathbf{x})$} (4.0, -1.8);
    \draw[thick, -{Stealth[length=2.5mm]}, analogygreen] (5.6, -1.8) -- node[below, font=\scriptsize\sffamily] {Reconstruction: $g_\phi(\mathbf{z})$} (8.8, -1.8);
\end{tikzpicture}
\caption{Undercomplete Autoencoder architecture. An input vector $\mathbf{x} \in \mathbb{R}^D$ is compressed by encoder $f_\theta$ into low-dimensional bottleneck $\mathbf{z} \in \mathbb{R}^d$ ($d \ll D$). Decoder $g_\phi$ reconstructs $\hat{\mathbf{x}}$. By training exclusively on healthy machinery telemetry, the network learns to reconstruct nominal patterns accurately, but fails to reconstruct mechanical faults.}\label{fig:autoencoder_arch}
\end{figure}

The network consists of two functional stages:
\begin{enumerate}
    \item \textbf{Encoder} $f_\theta$: Maps the input vector $\mathbf{x}$ to latent code $\mathbf{z}$:
    \begin{equation}
    \mathbf{z} = f_\theta(\mathbf{x}) = \sigma(\mathbf{W}_e \mathbf{x} + \mathbf{b}_e)\label{eq:encoder}
    \end{equation}
    \item \textbf{Decoder} $g_\phi$: Maps latent code $\mathbf{z}$ back to input space:
    \begin{equation}
    \hat{\mathbf{x}} = g_\phi(\mathbf{z}) = \sigma(\mathbf{W}_d \mathbf{z} + \mathbf{b}_d)\label{eq:decoder}
    \end{equation}
\end{enumerate}
During training on clean, healthy operational baseline data $\mathcal{D}_{\text{healthy}}$, network parameters $(\theta, \phi)$ are optimized to minimize the Mean Squared Error (MSE) reconstruction loss:
\begin{equation}
\mathcal{L}_{\text{MSE}}(\theta, \phi) = \frac{1}{N}\sum_{i=1}^N \|\mathbf{x}_i - g_\phi(f_\theta(\mathbf{x}_i))\|^2 = \frac{1}{N}\sum_{i=1}^N \frac{1}{D}\sum_{j=1}^D (x_{i,j} - \hat{x}_{i,j})^2\label{eq:autoencoder_loss}
\end{equation}

\subsubsection{The Reconstruction Error Anomaly Score}

Because the network was trained exclusively on healthy patterns, its weights $\mathbf{W}_e, \mathbf{W}_d$ encode the non-linear physical manifold of nominal operation.

When a normal telemetry vector $\mathbf{x}_{\text{normal}}$ is passed through the network, the autoencoder compresses it into the bottleneck and reconstructs it with high fidelity:
\begin{equation*}
e(\mathbf{x}_{\text{normal}}) = \|\mathbf{x}_{\text{normal}} - \hat{\mathbf{x}}_{\text{normal}}\|^2 \approx 0
\end{equation*}
When an anomalous vector $\mathbf{x}_{\text{fault}}$ arrives (such as a bearing spalling resonance or cracked gear tooth vibration burst), the bottleneck representation cannot preserve features it has never observed. The decoder produces a flawed reconstruction, resulting in a large \textbf{Reconstruction Error Spike}:
\begin{equation}
\text{Reconstruction Error } e(\mathbf{x}) = \frac{1}{D}\sum_{j=1}^D (x_j - \hat{x}_j)^2\label{eq:recon_error}
\end{equation}

\begin{figure}[htbp]
\centering
\begin{tikzpicture}[scale=0.92]
    % Time axis
    \draw[-{Stealth[length=2.5mm]}, thick, gray!80] (0, 0) -- (11.5, 0) node[right, font=\small\sffamily] {Time $t$};
    \draw[-{Stealth[length=2.5mm]}, thick, gray!80] (0, 0) -- (0, 4.2) node[above, font=\small\sffamily] {Reconstruction Error $e(t)$};

    % Threshold Line
    \draw[warningred, dashed, thick] (0, 2.2) -- (11.0, 2.2) node[right, font=\footnotesize\bfseries\sffamily] {Threshold $\tau = 0.25$};

    % Healthy trace (t = 0 to 6)
    \draw[keyideablue, thick] plot[smooth] coordinates {
        (0.2, 0.4) (1.0, 0.6) (1.8, 0.35) (2.5, 0.7) (3.2, 0.45) 
        (4.0, 0.55) (4.8, 0.4) (5.5, 0.65) (6.0, 0.5)
    };
    \node[keyideablue, font=\footnotesize\sffamily] at (3.0, 1.0) {\textbf{Healthy Operation} ($e < \tau$)};

    % Fault onset (t = 6 to 10.5)
    \draw[warningred, very thick] plot[smooth] coordinates {
        (6.0, 0.5) (6.5, 1.2) (7.0, 2.0) (7.5, 2.9) (8.2, 3.6) 
        (9.0, 3.2) (9.8, 3.8) (10.5, 3.5)
    };
    \fill[warningred] (7.2, 2.2) circle (3pt);
    \draw[-{Stealth[length=2mm]}, warningred, thick] (7.2, 1.2) -- (7.2, 2.0);
    \node[warningred, font=\scriptsize\bfseries\sffamily, anchor=north] at (7.2, 1.1) {Fault Trigger Event};

    \node[warningred, font=\footnotesize\sffamily] at (9.0, 2.8) {\textbf{Alarm Condition} ($e > \tau$)};
\end{tikzpicture}
\caption{Reconstruction error time series. During healthy operation, the autoencoder reconstructs inputs with low residual error. As physical wear degrades the mechanism, reconstruction error surges past threshold $\tau$, triggering an alert.}\label{fig:recon_error_timeseries}
\end{figure}

\paragraph{Selecting the Anomaly Threshold $\tau$}
Without labeled anomalies, the threshold $\tau$ is established using a held-out clean validation set $\mathcal{D}_{\text{val}}$:
\begin{itemize}
    \item \textbf{Percentile Rule}: Set $\tau = \text{Percentile}(e(\mathcal{D}_{\text{val}}), 99.5\%)$, fixing the nominal false alarm rate to exactly $0.5\%$.
    \item \textbf{Extreme Value Theory (EVT)}: Model the tail of the reconstruction error distribution using the Generalized Pareto Distribution (GPD) to establish statistically rigorous extreme threshold bounds without Gaussian assumptions.
\end{itemize}

% ------------------------------------------------------------------------------
\subsection{Comparative Case Study: Industrial Pumping Station}\label{subsec:comparative_case_study}
% ------------------------------------------------------------------------------

To empirically evaluate the four anomaly detection methods under realistic conditions, we implement an end-to-end benchmark on synthetic telemetry from an industrial pumping station.

The system monitors three coupled physical features:
\begin{enumerate}
    \item Motor Temperature ($x_1$, nominal $30.0^\circ\text{C}$, $\sigma = 2.0^\circ\text{C}$)
    \item Vibration RMS ($x_2$, nominal $0.20\text{ mm/s}$, $\sigma = 0.07\text{ mm/s}$)
    \item Discharge Pressure ($x_3$, nominal $5.0\text{ bar}$, $\sigma = 0.5\text{ bar}$)
\end{enumerate}

The test set contains $N_{\text{total}} = 525$ samples ($500$ nominal operating cycles and $25$ ground-truth physical faults):
\begin{itemize}
    \item $10$ severe point spikes (extreme overheating or excessive mechanical vibration).
    \item $15$ correlation-breaking faults (individual sensor values reside within nominal ranges, but physical correlation is violated, e.g., elevated temperature accompanied by sudden pressure drop).
\end{itemize}

\begin{lstlisting}[style=pythonstyle, caption={End-to-end comparative anomaly detection benchmark.}, label={lst:anomaly_benchmark}]
import numpy as np
import pandas as pd
from sklearn.preprocessing import StandardScaler
from sklearn.ensemble import IsolationForest
from sklearn.svm import OneClassSVM
from sklearn.decomposition import PCA
from scipy.spatial.distance import mahalanobis
from scipy.stats import chi2
from sklearn.metrics import confusion_matrix

np.random.seed(42)

# 1. Synthesize baseline telemetry: 500 normal samples, 25 anomalies
n_normal, n_anom = 500, 25
mean_normal = [30.0, 0.20, 5.0]
cov_normal = [
    [4.0,  0.08, 0.6],
    [0.08, 0.005, 0.02],
    [0.6,  0.02, 0.25]
]

# Generate healthy operating distribution
X_normal = np.random.multivariate_normal(mean_normal, cov_normal, n_normal)
y_normal = np.zeros(n_normal, dtype=int)

# Generate anomalies: 10 point spikes + 15 correlation-breaking faults
anom_spikes = np.random.multivariate_normal([45.0, 0.60, 5.0], cov_normal, 10)
anom_corr = np.column_stack([
    np.random.normal(38.0, 1.0, 15),
    np.random.normal(0.20, 0.02, 15),
    np.random.normal(3.2, 0.2, 15)
])
X_anom = np.vstack([anom_spikes, anom_corr])
y_anom = np.ones(n_anom, dtype=int)

X = np.vstack([X_normal, X_anom])
y = np.concatenate([y_normal, y_anom])

# Standardize features
scaler = StandardScaler()
X_scaled = scaler.fit_transform(X)

# Model 1: Mahalanobis Distance (Fitted on normal baseline)
mu_hat = np.mean(X_scaled[:n_normal], axis=0)
cov_hat = np.cov(X_scaled[:n_normal], rowvar=False)
inv_cov = np.linalg.inv(cov_hat)
m_dists = np.array([mahalanobis(x, mu_hat, inv_cov) for x in X_scaled])
thresh_m = np.sqrt(chi2.ppf(0.99, df=3)) # 99% Chi-square critical value
pred_m = (m_dists > thresh_m).astype(int)

# Model 2: Isolation Forest (contamination = 25 / 525 ~ 0.048)
iso = IsolationForest(contamination=0.048, random_state=42)
pred_iso = (iso.fit_predict(X_scaled) == -1).astype(int)

# Model 3: One-Class SVM (nu = 0.048, RBF kernel)
oc_svm = OneClassSVM(nu=0.048, kernel='rbf', gamma='scale')
pred_svm = (oc_svm.fit_predict(X_scaled) == -1).astype(int)

# Model 4: PCA Reconstruction Error (k = 2 components)
pca = PCA(n_components=2)
X_pca = pca.fit_transform(X_scaled)
X_recon = pca.inverse_transform(X_pca)
recon_err = np.mean((X_scaled - X_recon)**2, axis=1)
thresh_pca = np.percentile(recon_err[:n_normal], 95)
pred_pca = (recon_err > thresh_pca).astype(int)

# Benchmark Output Table
models = [('Mahalanobis Distance', pred_m), 
          ('Isolation Forest', pred_iso), 
          ('One-Class SVM', pred_svm), 
          ('PCA Reconstruction', pred_pca)]

print(f"{'Algorithm':22s} | {'Caught (TP)':11s} | {'False Alarm (FP)':16s} | {'Precision':9s} | {'Recall':6s} | {'F1-Score':8s}")
print("-" * 85)
for name, pred in models:
    cm = confusion_matrix(y, pred)
    tp, fp, fn = cm[1,1], cm[0,1], cm[1,0]
    rec = tp / (tp + fn)
    prec = tp / (tp + fp) if (tp + fp) > 0 else 0
    f1 = 2 * prec * rec / (prec + rec) if (prec + rec) > 0 else 0
    print(f"{name:22s} | {tp:2d}/25       | {fp:2d}               | {prec:.2f}      | {rec:.2f}   | {f1:.3f}")
\end{lstlisting}

\begin{lstlisting}[style=outputstyle, caption={Benchmark execution output on pump fault dataset.}]
Algorithm              | Caught (TP) | False Alarm (FP) | Precision | Recall | F1-Score
-------------------------------------------------------------------------------------
Mahalanobis Distance   | 25/25       |  2               | 0.93      | 1.00   | 0.962   
Isolation Forest       | 20/25       |  6               | 0.77      | 0.80   | 0.784   
One-Class SVM          |  9/25       | 18               | 0.33      | 0.36   | 0.346   
PCA Reconstruction     |  1/25       | 25               | 0.04      | 0.04   | 0.039   
\end{lstlisting}

\subsubsection{Benchmark Analysis and Method Selection Guide}

Table~\ref{tab:anomaly_algorithm_comparison} synthesizes the architectural trade-offs across all four paradigms.

\begin{table}[htbp]
\centering
\small
\begin{tabular}{lp{3.5cm}p{3.5cm}p{2.6cm}c}
\toprule
\textbf{Algorithm} & \textbf{Core Strengths} & \textbf{Key Limitations} & \textbf{Complexity (Inference)} & \textbf{MCU Ready?} \\
\midrule
\textbf{Z-Score / Tukey} & Instantaneous; zero memory; explainable. & Blind to feature correlations. & $\mathcal{O}(D)$ & Yes (8-bit) \\
\addlinespace[3pt]
\textbf{Mahalanobis} & Catches linear correlation breaks; optimal for Gaussian data. & Fails on multi-modal distributions; requires matrix inversion. & $\mathcal{O}(D^2)$ & Yes (ESP32) \\
\addlinespace[3pt]
\textbf{Isolation Forest} & Robust to contamination; no distance metric needed; fast. & Axis-aligned splits struggle with complex non-linear manifolds. & $\mathcal{O}(t \cdot \log \psi)$ & Yes (C-arrays) \\
\addlinespace[3pt]
\textbf{One-Class SVM} & Flexible non-linear boundary in kernel space. & Quadratic scaling $\mathcal{O}(N^2)$; hypersensitive to $\gamma$. & $\mathcal{O}(N_{\text{SV}} \cdot D)$ & Borderline \\
\addlinespace[3pt]
\textbf{Autoencoder} & Unrivaled on complex non-linear high-dim spectra. & Requires large normal training set; opaque latent space. & $\mathcal{O}(\text{FLOPs})$ & Yes (TFLite Micro) \\
\bottomrule
\end{tabular}
\caption{Comprehensive architectural comparison of anomaly detection algorithms for IoT and edge deployment.}\label{tab:anomaly_algorithm_comparison}
\end{table}

\begin{iotbox}[Deploying Anomaly Detectors on Edge Microcontrollers (ESP32)]
Why run anomaly detection directly on the edge instead of streaming raw telemetry to AWS or Azure?
\begin{itemize}
    \item \textbf{Sub-Millisecond E-Stop}: When a mechanical spindle bearing seizes at $24{,}000\text{ RPM}$, severe mechanical damage occurs within $50\text{ milliseconds}$. Cellular/LoRaWAN network latency ($200\text{--}2{,}000\text{ ms}$) is far too slow to prevent machine destruction.
    \item \textbf{Bandwidth Savings}: A 3-axis accelerometer sampled at $10\text{ kHz}$ generates $60\text{ kB/second}$ ($5.18\text{ GB/day}$). An ESP32 computing on-device Mahalanobis distance or TFLite autoencoder reconstruction reduces this continuous firehose down to a 1-byte status flag transmitted only upon fault onset.
    \item \textbf{C Implementation of Mahalanobis}: The mean vector $\boldsymbol{\mu} \in \mathbb{R}^D$ and precomputed inverse covariance matrix $\mathbf{\Sigma}^{-1} \in \mathbb{R}^{D \times D}$ require only $D + D^2$ 32-bit floats. For $D = 4$, this occupies merely $80\text{ bytes}$ of flash memory! Evaluating $D_M^2$ requires a $4 \times 4$ vector-matrix multiplication, completing in under $2\text{ microseconds}$ on a $240\text{ MHz}$ ESP32.
\end{itemize}
\end{iotbox}

% ------------------------------------------------------------------------------
\subsection{Summary and Practical Checklist}\label{subsec:anomaly_summary}
% ------------------------------------------------------------------------------

\begin{takeaways}[Essential Anomaly Detection Principles]
\begin{itemize}
    \item \textbf{Anomalies are Rare and Novel}: Never treat fault detection as a standard balanced classification problem. Supervised classifiers fail when confronted with unprecedented failure modes.
    \item \textbf{Univariate Bounds Miss Physical Breaks}: Independent thresholds fail whenever variables are physically coupled. Always deploy multivariate methods (Mahalanobis distance, Isolation Forest, Autoencoders) to monitor interdependent sensors.
    \item \textbf{Subsampling is Crucial in iForest}: Restricting tree construction to sub-samples of size $\psi = 256$ neutralizes swamping and masking while dramatically speeding up training to $\mathcal{O}(t \cdot \psi \log \psi)$.
    \item \textbf{The Meaning of $\nu$ in OC-SVM}: The hyperparameter $\nu$ simultaneously sets an upper limit on training false alarms and a lower limit on the fraction of support vectors.
    \item \textbf{Autoencoders Profile Normalcy}: By constraining capacity through an informational bottleneck, autoencoders learn the manifold of healthy operation; unseen anomalies trigger immediate reconstruction error spikes.
    \item \textbf{Edge Deployment Saves Machines}: Computing anomaly scores directly on embedded microcontrollers enables sub-millisecond emergency shutdowns while reducing transmission bandwidth by over $99.9\%$.
\end{itemize}
\end{takeaways}

% ------------------------------------------------------------------------------
\subsection{Check Yourself}\label{subsec:anomaly_check_yourself}
% ------------------------------------------------------------------------------

\begin{checkyourself}[Practice Exercises]
\begin{enumerate}
    \item \textbf{Mahalanobis Distance and Correlation}:
    A smart pump sensor tracks motor temperature $x_1$ ($^\circ\text{C}$) and vibration $x_2$ ($\text{mm/s}$). The healthy baseline parameters are:
    \begin{equation*}
    \boldsymbol{\mu} = \begin{bmatrix} 30.0 \\ 0.20 \end{bmatrix}, \qquad \mathbf{\Sigma} = \begin{bmatrix} 4.0 & 0.04 \\ 0.04 & 0.0016 \end{bmatrix}
    \end{equation*}
    A new sensor observation arrives: $\mathbf{x} = [32.0, \ 0.24]^T$.
    \begin{enumerate}[label=(\alph*)]
        \item Compute the difference vector $\mathbf{d} = \mathbf{x} - \boldsymbol{\mu}$.
        \item Calculate $\det(\mathbf{\Sigma})$ and derive the inverse covariance matrix $\mathbf{\Sigma}^{-1}$.
        \item Compute the squared Mahalanobis distance $D_M^2(\mathbf{x})$.
        \item Given the Chi-squared critical threshold $\chi^2_{2, 0.95} = 5.991$, is this point anomalous?
    \end{enumerate}

    \item \textbf{Isolation Forest Path Length and Anomaly Score}:
    An Isolation Forest is trained with subsampling size $\psi = 256$.
    \begin{enumerate}[label=(\alph*)]
        \item Calculate the average BST path length normalizer $c(256)$ using $\gamma \approx 0.5772$.
        \item Observation $\mathbf{x}_A$ has an average path length $\mathbb{E}[h(\mathbf{x}_A)] = 4.2$. Compute its anomaly score $s(\mathbf{x}_A, 256)$.
        \item Observation $\mathbf{x}_B$ has an average path length $\mathbb{E}[h(\mathbf{x}_B)] = 11.5$. Compute its anomaly score $s(\mathbf{x}_B, 256)$.
        \item Using decision threshold $s \ge 0.60$, classify both points.
    \end{enumerate}

    \item \textbf{Tukey's IQR Fences for Robust Temperature Monitoring}:
    An environmental sensor records greenhouse temperatures over a 24-hour period. Ordered quartile analysis yields $Q_1 = 22.0^\circ\text{C}$ and $Q_3 = 28.0^\circ\text{C}$.
    \begin{enumerate}[label=(\alph*)]
        \item Compute the Interquartile Range ($\text{IQR}$).
        \item Determine the inner lower and upper fences ($F_{\text{lower}}, F_{\text{upper}}$).
        \item Classify two new readings: $x_1 = 11.5^\circ\text{C}$ and $x_2 = 35.0^\circ\text{C}$.
    \end{enumerate}

    \item \textbf{Autoencoder Reconstruction Error Computation}:
    A 4-dimensional sensor vector $\mathbf{x} = [1.2, -0.8, 3.5, 0.0]^T$ is passed through a trained autoencoder, producing reconstructed output $\hat{\mathbf{x}} = [1.1, -0.7, 2.1, 0.2]^T$.
    \begin{enumerate}[label=(\alph*)]
        \item Compute the error residual vector $\mathbf{e} = \mathbf{x} - \hat{\mathbf{x}}$.
        \item Calculate the Mean Squared Error (MSE) reconstruction score.
        \item If the calibrated normal threshold is $\tau = 0.25$, is this sample anomalous? Which specific sensor channel was primarily responsible?
    \end{enumerate}

    \item \textbf{IoT Anomaly Classification in Practice}:
    In an edge IoT facility, identify whether each of the following scenarios represents a \textit{Point Anomaly}, a \textit{Contextual Anomaly}, or a \textit{Collective Anomaly}:
    \begin{enumerate}[label=(\alph*)]
        \item A single temperature reading on a furnace flashes to $999^\circ\text{C}$ for one sample due to an ADC SPI bus glitch.
        \item An office building HVAC system draws $45\text{ kW}$ of electrical power continuously at 3:00 AM on Sunday morning.
        \item A vibration sensor on an electric motor exhibits a regular sinusoidal beat frequency pattern over a 30-second window, even though every individual sample amplitude remains within nominal $3\sigma$ bounds.
    \end{enumerate}
\end{enumerate}
\end{checkyourself}

% ------------------------------------------------------------------------------
\subsection*{Answers}\label{subsec:anomaly_answers}
% ------------------------------------------------------------------------------

\begin{enumerate}
    \item \textbf{Mahalanobis Distance and Correlation}:
    \begin{enumerate}[label=(\alph*)]
        \item \textbf{Difference Vector}:
        \begin{equation*}
        \mathbf{d} = \mathbf{x} - \boldsymbol{\mu} = \begin{bmatrix} 32.0 - 30.0 \\ 0.24 - 0.20 \end{bmatrix} = \begin{bmatrix} 2.0 \\ 0.04 \end{bmatrix}
        \end{equation*}

        \item \textbf{Covariance Determinant and Matrix Inversion}:
        \begin{equation*}
        \det(\mathbf{\Sigma}) = (4.0)(0.0016) - (0.04)^2 = 0.0064 - 0.0016 = \mathbf{0.0048}
        \end{equation*}
        Using the analytical $2 \times 2$ inverse formula $\mathbf{\Sigma}^{-1} = \frac{1}{\det(\mathbf{\Sigma})} \begin{bmatrix} \sigma_{22} & -\sigma_{12} \\ -\sigma_{21} & \sigma_{11} \end{bmatrix}$:
        \begin{equation*}
        \mathbf{\Sigma}^{-1} = \frac{1}{0.0048} \begin{bmatrix} 0.0016 & -0.04 \\ -0.04 & 4.0 \end{bmatrix} = \begin{bmatrix} \frac{1}{3} & -\frac{25}{3} \\ -\frac{25}{3} & \frac{2500}{3} \end{bmatrix} \approx \begin{bmatrix} 0.3333 & -8.3333 \\ -8.3333 & 833.3333 \end{bmatrix}
        \end{equation*}

        \item \textbf{Squared Mahalanobis Distance}:
        \begin{align*}
        \mathbf{\Sigma}^{-1} \mathbf{d} &= \begin{bmatrix} \frac{1}{3}(2.0) - \frac{25}{3}(0.04) \\ -\frac{25}{3}(2.0) + \frac{2500}{3}(0.04) \end{bmatrix} = \begin{bmatrix} \frac{2.0 - 1.0}{3} \\ \frac{-50.0 + 100.0}{3} \end{bmatrix} = \begin{bmatrix} \frac{1}{3} \\ \frac{50}{3} \end{bmatrix} \approx \begin{bmatrix} 0.3333 \\ 16.6667 \end{bmatrix} \\
        D_M^2(\mathbf{x}) &= \mathbf{d}^T (\mathbf{\Sigma}^{-1} \mathbf{d}) = 2.0\left(\frac{1}{3}\right) + 0.04\left(\frac{50}{3}\right) = \frac{2.0}{3} + \frac{2.0}{3} = \frac{4.0}{3} \approx \mathbf{1.3333}
        \end{align*}
        $D_M(\mathbf{x}) = \sqrt{1.3333} \approx 1.1547$.

        \item \textbf{Anomaly Decision}:
        Because $D_M^2(\mathbf{x}) = 1.3333 < \chi^2_{2, 0.95} = 5.991$, the observation falls safely within the $95\%$ confidence ellipse. The point is classified as \textbf{Normal}.
    \end{enumerate}

    \item \textbf{Isolation Forest Path Length and Anomaly Score}:
    \begin{enumerate}[label=(\alph*)]
        \item \textbf{BST Average Path Length Normalizer $c(256)$}:
        \begin{align*}
        c(256) &= 2 \ln(256 - 1) + 2(0.57721566) - \frac{2(256 - 1)}{256} \\
        &= 2 \ln(255) + 1.154431 - \frac{510}{256} \\
        &= 2(5.5412635) + 1.154431 - 1.9921875 \\
        &= 11.082527 + 1.154431 - 1.9921875 = \mathbf{10.2448}
        \end{align*}

        \item \textbf{Anomaly Score for $\mathbf{x}_A$ ($\mathbb{E}[h] = 4.2$)}:
        \begin{equation*}
        s(\mathbf{x}_A, 256) = 2^{-\frac{4.2}{10.2448}} = 2^{-0.40996} \approx \mathbf{0.7526}
        \end{equation*}

        \item \textbf{Anomaly Score for $\mathbf{x}_B$ ($\mathbb{E}[h] = 11.5$)}:
        \begin{equation*}
        s(\mathbf{x}_B, 256) = 2^{-\frac{11.5}{10.2448}} = 2^{-1.1225} \approx \mathbf{0.4593}
        \end{equation*}

        \item \textbf{Classification ($s \ge 0.60$)}:
        \begin{itemize}
            \item Point $\mathbf{x}_A$: $s = 0.7526 \ge 0.60 \implies$ \textbf{Anomalous Outlier} (isolated very quickly in shallow leaf nodes).
            \item Point $\mathbf{x}_B$: $s = 0.4593 < 0.60 \implies$ \textbf{Normal Instance} (deeply nested in dense cluster, requiring more splits than average).
        \end{itemize}
    \end{enumerate}

    \item \textbf{Tukey's IQR Fences for Robust Temperature Monitoring}:
    \begin{enumerate}[label=(\alph*)]
        \item \textbf{Interquartile Range}:
        \begin{equation*}
        \text{IQR} = Q_3 - Q_1 = 28.0^\circ\text{C} - 22.0^\circ\text{C} = \mathbf{6.0^\circ\text{C}}
        \end{equation*}

        \item \textbf{Inner Fences}:
        \begin{align*}
        F_{\text{lower}} &= Q_1 - 1.5 \times \text{IQR} = 22.0 - 1.5(6.0) = 22.0 - 9.0 = \mathbf{13.0^\circ\text{C}} \\
        F_{\text{upper}} &= Q_3 + 1.5 \times \text{IQR} = 28.0 + 1.5(6.0) = 28.0 + 9.0 = \mathbf{37.0^\circ\text{C}}
        \end{align*}

        \item \textbf{Classification}:
        \begin{itemize}
            \item $x_1 = 11.5^\circ\text{C} < 13.0^\circ\text{C} \implies$ \textbf{Anomalous Outlier} (violates lower fence, indicating possible sensor disconnect or unheated freezing condition).
            \item $x_2 = 35.0^\circ\text{C} \in [13.0^\circ\text{C}, \ 37.0^\circ\text{C}] \implies$ \textbf{Normal} (within allowable inner fences).
        \end{itemize}
    \end{enumerate}

    \item \textbf{Autoencoder Reconstruction Error Computation}:
    \begin{enumerate}[label=(\alph*)]
        \item \textbf{Error Residual Vector}:
        \begin{equation*}
        \mathbf{e} = \mathbf{x} - \hat{\mathbf{x}} = \begin{bmatrix} 1.2 - 1.1 \\ -0.8 - (-0.7) \\ 3.5 - 2.1 \\ 0.0 - 0.2 \end{bmatrix} = \begin{bmatrix} 0.1 \\ -0.1 \\ 1.4 \\ -0.2 \end{bmatrix}
        \end{equation*}

        \item \textbf{Mean Squared Error (MSE)}:
        \begin{align*}
        \|\mathbf{e}\|^2 &= (0.1)^2 + (-0.1)^2 + (1.4)^2 + (-0.2)^2 = 0.01 + 0.01 + 1.96 + 0.04 = \mathbf{2.02} \\
        \text{MSE} &= \frac{1}{D} \|\mathbf{e}\|^2 = \frac{2.02}{4} = \mathbf{0.5050}
        \end{align*}

        \item \textbf{Decision and Attribution}:
        Because $\text{MSE} = 0.5050 > \tau = 0.25$, the reading is flagged as an \textbf{Anomaly}. Looking at the squared residuals, Sensor Channel 3 accounts for $(1.4)^2 / 2.02 = 1.96 / 2.02 \approx \mathbf{97.0\%}$ of the total error, directly isolating the fault to Sensor 3.
    \end{enumerate}

    \item \textbf{IoT Anomaly Classification in Practice}:
    \begin{enumerate}[label=(\alph*)]
        \item \textbf{Point Anomaly}: An instantaneous single-sample outlier ($999^\circ\text{C}$) that deviates completely from the allowable physical range due to hardware communication corruption.
        \item \textbf{Contextual Anomaly}: A continuous power draw of $45\text{ kW}$ is entirely normal during peak manufacturing hours on Wednesday afternoon, but is anomalous in the context of Sunday at 3:00 AM when all machinery should be powered down.
        \item \textbf{Collective Anomaly}: Every individual vibration reading sits inside the nominal $3\sigma$ amplitude envelope, but the collective 30-second temporal trajectory forms an abnormal oscillation pattern indicating mechanical resonance or bearing deterioration.
    \end{enumerate}
\end{enumerate}
